\section{Конструктивное исследование и перенос}\label{sec:investigation}
Достаточность задаёт критерий для фиксированного наблюдения, но не стратегию приобретения информации. Сначала синтезируем конечные адаптивные исследования, затем выясним, когда синтезированную стратегию можно перенести на другую задачу. Объявленные списки миров, действий и портов конечны; равенство портов разрешимо. Оценка эффективности исчерпывающего синтеза не утверждается.

\subsection{Синтез адаптивного исследования}
Задача задаётся списком миров $E$, непустым списком действий $\mathcal A$, списком портов $Q$, функцией чтения $r(q,w)\in\{0,1\}$, ценами портов $c(q)\in\Nat$, терминальными потерями $L(a,w)\in\Nat$ и пределом глубины $d$. Исследование --- конечное дерево
\[
 t ::= \operatorname{act}(a)\mid\operatorname{observe}(q,t_0,t_1).
\]
Допустимость требует объявленных действий и портов, глубины не более $d$ и отсутствия повторного порта вдоль пути (все вхождения приобретённого порта удаляются из оставшегося списка). Для реализованной стоимости $J_w(t)$ семантика такова: $J_w(\operatorname{act}(a))=L(a,w)$ и $J_w(\operatorname{observe}(q,t_0,t_1))=c(q)+J_w(t_{r(q,w)})$. Оплачивается только реализованная ветвь. Пусть $W_E(t)=\max_{w\in E}J_w(t)$, максимум пустого списка равен нулю, а $V_d(E,Q)$ --- оптимальная худшая стоимость допустимого дерева. Обозначим через $E_b(q)$ фильтрацию $E$ по значению чтения $b$, и $L_*(E)=\min_{a\in\mathcal A}\max_{w\in E}L(a,w)$.

\begin{theorem}[Оптимальное конечное исследование]\label{thm:investigation}
Исполнимый решатель достигает $V_d(E,Q)$ допустимым деревом. При $E\ne\varnothing$,
\begin{align}
 V_0(E,Q)&=L_*(E),\\
 V_{d+1}(E,Q)&=\min\Bigl(\{L_*(E)\}\cup\{c(q)+\max_{b\in\{0,1\}} \\
 &\hspace{24mm} V_d(E_b(q),Q\setminus\{q\}):q\in Q\}\Bigr).
\end{align}
Для пустого $E$ значение равно нулю. Для жёсткого допуска терминальных потерь $\epsilon$ и бюджета наблюдений $b$ заменим терминальную потерю на $0$, если $L(a,w)\le\epsilon$, и на $b+1$ иначе. Оптимум такой ограниченной модели не превосходит $b$ тогда и только тогда, когда допустимое дерево соблюдает обе границы во всех объявленных мирах.
\end{theorem}
\begin{proof}
Индукция по глубине. Любое допустимое дерево либо является терминальным действием, либо начинается с объявленного порта и двух допустимых поддеревьев меньшей глубины. Для непустого $E$ худшая стоимость узла наблюдения равна плате плюс максимум стоимости ветвей: одна из ветвей реализуется, и все стоимости неотрицательны. Решатель выбирает наименее дорогого кандидата среди оптимального терминального действия и рекурсивно синтезированных узлов наблюдения. По индукции стоимость любого конкурирующего дерева не меньше минимума, а выбранный кандидат допустим. В ограниченной модели нарушающий допуск лист стоит не менее $b+1$, допустимый даёт нуль; поэтому ограничение худшей стоимости в точности проверяет жёсткую достижимость. См. \rc{solve\_legal}, \rc{solve\_optimal}, \rc{bellman} и \rc{epsilon\_budget\_iff} в \rc{AdaptiveObservation.lean}.
\end{proof}
Адаптивный оптимум не превосходит лучшего допустимого фиксированного расписания (\rc{adaptive\_le\_static}). В свидетеле-мультиплексоре читается бит селектора, затем указанный бит данных: стоимость $2$ вместо статического оптимума $3$ при единичных ценах и потере $4$ за неверный ответ. Свидетель XOR также исключает общее жадное правило одношаговой ценности. Синтез точен в объявленном языке деревьев, а не по произвольным программам.

\subsection{Перенос и экономика библиотеки}
Для задач $S,T$ морфизм переноса $M=(h,p,a)$ отображает целевые миры в исходные, исходные порты в целевые, исходные действия в целевые, сохраняя объявленное членство. Карта портов инъективна на объявленных исходных портах, $d_S\le d_T$, и для всех объявленных аргументов
\begin{align*}
 r_T(p(q),w)&=r_S(q,h(w)), & c_T(p(q))&\le c_S(q),\\
 L_T(a(u),w)&\le L_S(u,h(w)).
\end{align*}
Перенос $M_*t$ --- синтаксическая подстановка меток портов и действий, не зависящая от неизвестного целевого мира.
\begin{theorem}[Сертифицированный перенос]\label{thm:transfer}
Для любого допустимого в исходной задаче дерева $t$ дерево $M_*t$ допустимо в целевой задаче, и $J^T_w(M_*t)\le J^S_{h(w)}(t)$ для каждого объявленного целевого мира. Плата за наблюдения и терминальные потери не возрастают по отдельности. Следовательно, $V_T\le V_S$; взаимные морфизмы влекут равенство оптимумов. Тождественность и композиция морфизмов индуцируют тождественность и композицию переносов деревьев.
\end{theorem}
\begin{proof}
Структурная индукция по дереву даёт совпадение выбора ветвей по равенству чтений и неравенства стоимости по сложению. Инъективность гарантирует, что удаление наблюдённого порта не удалит образ другого исходного порта. Неравенство глубин обеспечивает допустимость. Максимизация по целевым мирам и применение к исходному оптимизатору доказывают неравенство значений. Композиция есть композиция функций (контравариантная по мирам), а подстановки деревьев компонуются. Это результаты \rc{transfer\_legal}, \rc{transfer\_runCost\_le}, \rc{transfer\_hard\_contract}, \rc{optimum\_le}, \rc{mutual\_optimum\_eq} и \rc{transfer\_compose} в \rc{InvestigationTransfer.lean}.
\end{proof}
Равенство оптимумов само по себе не влечёт существования морфизма; конечные свидетели равных оптимумов демонстрируют нарушение инъективности и сохранения потерь. Перенос также требует затрат. Для указанного маршрута с накладными расходами $k$ и стоимостью целевого исполнения $e\ge V_T$ повторное использование строго дешевле нового синтеза за $D$ с оптимальным исполнением тогда и только тогда, когда $k+(e-V_T)<D$. Явный селектор библиотеки минимизирует стоимость по новому синтезу и перечисленным сертифицированным предложениям; при цене поиска $s$ и цене нового синтеза с исполнением $F$ поиск окупается тогда и только тогда, когда некоторое найденное предложение удовлетворяет $s+k+e<F$. Полнота относительно семантического оракула предложений требует полного по стоимости допущенного списка, а не просто корректного подмножества. Это \rc{paid\_candidate\_iff}, \rc{search\_profitable\_iff} и \rc{complete\_search\_matches\_oracle} в \rc{InvestigationLibrary.lean}.

\subsection{Конструктивное обнаружение морфизмов}
Теорема переноса полезна, лишь когда доступен её свидетель. Для индексированных конечных задач (объявлены все действия, включая действие по умолчанию) перечислим все карты $p:Q_S\to Q_T$ и $a:\mathcal A_S\to\mathcal A_T$. Для фиксированной пары определим
\[
 \mathcal Y_w(p,a)=\{y\in X_S: \forall q\in Q_S,\ r_T(p(q),w)=r_S(q,y),\quad
 \forall u\in\mathcal A_S,\ L_T(a(u),w)\le L_S(u,y)\}.
\]
\begin{theorem}[Каноническое обнаружение]\label{thm:discovery}
Фиксированная пара $(p,a)$ продолжается до морфизма тогда и только тогда, когда соблюдены условия глубины и портов и $\mathcal Y_w(p,a)\ne\varnothing$ для каждого целевого мира $w$. Выбор первого свидетеля в каждой конечной строке исходных миров строит морфизм. Список $\CanonRef=\rc{discoverAll}(S,T)$ содержит представителя с теми же картами портов и действий для каждого допустимого морфизма; он пуст тогда и только тогда, когда морфизмов нет.
\end{theorem}
\begin{proof}
Допустимая карта миров даёт элемент каждой строки. Обратно, каждому целевому миру можно независимо назначить первый элемент строки: морфизм не накладывает межмировых ограничений. Тем самым выполнены все обязательства чтения и потерь; отдельно проверенные условия глубины, цены и инъективности не меняются. Полный перебор карт включает пару любого допустимого морфизма. Непустой результат сам содержит элемент типа морфизмов. См. \rc{fixed\_pair\_iff}, \rc{discover\_complete} и \rc{empty\_iff\_no\_morphism} в \rc{MorphismDiscovery.lean}.
\end{proof}
Каноническому списку не требуется перечислять все возможные карты миров: перенесённое дерево зависит только от $p,a$. Это устраняет полный перебор глобальных карт миров, но не потенциально дорогой перебор остальных пар. Следующая задача --- сохранить эталонную семантику при прерывании, отсечении и обучении.
